\documentclass[UTF8,11pt,a4paper]{ctexart}
\usepackage[margin=23mm]{geometry}
\usepackage{amsmath,amssymb,booktabs,longtable,array}
\usepackage[colorlinks=true,urlcolor=blue,linkcolor=blue,citecolor=blue]{hyperref}
\usepackage{xurl}
\setlength{\emergencystretch}{3em}
\setlength{\tabcolsep}{4pt}
\title{18号赛题：Level目标的可达性与新阶段优化研究}
\author{项目独立实验记录}
\date{2026年10月11日；旧v14冻结报告另行保留}
\begin{document}
\maketitle
\begin{abstract}
本阶段首先复核官方正文和距离ISA原始论文，再实现逆向压力调度与严格共享根链调度。
官方正文仅列Level 0、1、3；Level 2、COPY占比分母和平均权重仍未定义。
不能将论文中的spill、重计算或硬件绕过当成赛题的COPY优化。
旧336图在原节点各执行一次、仅新增COPY的身份模型下，初始平均COPY/输出必要下界为6.690\%。
后续全DAG核心及共享归约证明将同336图必要下界提高至11.779\%，既有272合法输入的必要下界为10.086\%。
这排除了声明模型及这些合成输入上的等权10\%和2\%，不判定缺少官方数据、分母和权重的官方Level。
新60图持出中，共享根链候选全部正确，旧版本52图正确；共同52图COPY从33927降至18650。
同一批图的必要下界为10.772\%，实际11.052\%，因而即使大幅改善，也不能宣称达到10\%。
这些是保留全部失败的合成数据结论，不是官方Level通过声明。
\end{abstract}

\section{规范和合法性边界}
\href{https://www.chaspark.com/\#/races/competitions/1326690549003280384}{官方18题动态页}
在2026年10月11日重新读取，页面仍标最后更新2026年10月10日09:18。
问题是单basic block的DAG重排，GET为输入、SET为输出，引用\texttt{t\#d}指向当前指令前$d$个物理输出；
窗口$W=2^b$，$b\in\{2,3,4,5\}$，距离$W$本身合法。超距可插入COPY，必须保持功能正确。
Level 1要求额外COPY平均占比不超过10\%；Level 3要求不超过2\%并具有优秀时间复杂度。
正文没有Level 2，也没有说明按原指令数还是输出数计算、逐block还是按长度加权。
官网三指令GET/GET/MULW输入已作为回归；示意COPY输出按其相对引用定义会将两个操作数都指向第二个GET，独立验证器拒绝该输出。

默认模型保持每个原节点恰执行一次、运算符及有序操作数身份不变，只增加COPY。
GET/SET保留同寄存器RAW、WAR、WAW顺序，不同寄存器的访问允许重排；GET/GET没有额外顺序。
本地解释器的8位有符号、16位无符号、32位有符号整数回绕是测试配置，不冒称官方位宽。
官方未明确额外GET、spill、重计算、等值依赖替换、删除原节点是否允许。
已有等值别名策略只作为另一实验模型，不能复用原身份容量/COPY证书，也不进入本阶段默认成绩。
重复GET还可能涉及volatile、异常及读写时序；单凭读取同名寄存器不能授权复制读取。
即便将来允许额外GET，也应报告全部新增指令，不能利用只数COPY而扩大分母获得虚假改进。

\section{最直接相关的原始工作}
Sugita等CC2023论文\cite{sugita}给出STRAIGHT距离ISA的可靠代码生成。
已完整取得作者12页PDF并提取全文，证明和伪代码另渲染核实。
三个阶段是spill、SSA destruction和move insertion。其完备性依赖spill阶段限制活跃变量，不能推导任何COPY-only图都可行。
Lemma 1利用空闲窗口位置交换活跃值；Theorem 2给出活跃变量严格少于窗口的充分条件。
活跃数等于窗口并不因此一律无解。跨block、phi和调用约束超出本题单block范围。
最大距离接力的固定序列论证也不是允许全部重排后的全局COPY最优证明。
本项目允许当前原指令直接最后消费到期值，从而省去无必要接力。

Koizumi等ICCD2021\cite{koizumi}的距离调整覆盖跨block合流和反馈超边；aggressive spill会改变代码机制。
其真实程序指令数和模拟性能结果不是本赛题的COPY比例。
Sethi--Ullman\cite{sethi}解决树形表达式寄存器需求与求值顺序。
它解释先计算需求较高子树的价值，但自由寄存器树最优性不能转移至共享DAG及物理年龄窗口。

Chen的作者预印本v4\cite{chen}研究逆向min-reg调度、立即降低逆向活跃集的选择，以及共享消费者cluster，并使用CP-SAT作参照。
本阶段移植逆向就绪列表和局部解锁信息；保留所有数据和内存边，再交给距离发射器处理COPY。
后续实现单个持续DFS栈、至多64个目标的有限消费者后继聚类；仍非该文完整cluster或CP-SAT，
也没有把其GPU负载实验直接代入本题。
Register Bridging的ICCD2025机构记录\cite{bridge}研究硬件圆缓冲重定向绕过RMOV动态执行。
其动态绕过比例和模拟速度提升不会删掉静态COPY，不能证明软件达到2\%。

\section{独立下界与Level可达性}
若原身份$x$具有$u_x$个不同的原consumer，使用$c_x$次COPY接力，则出生到最后使用至少包含$u_x$次原消费和$c_x$次自身接力。
每段距离最多$W$，因此
\[
 u_x+c_x\le(c_x+1)W,\qquad
 c_x\ge\max\left(0,\left\lceil\frac{u_x-W}{W-1}\right\rceil\right).
\]
对身份求和是合法下界，因为每个COPY只属于一个原身份。
还有强制脊柱年龄和、双根阶段年龄和下界；逐图取已独立证明的最大值，未知图保留下界0。
原JSON检查器重新建立有序引用、独立consumer集合和适用结构，不调用调度器内部图。
证书只适用于明确的纯运算、唯一且读写分离的GET/SET寄存器，以及原身份单次、COPY-only变换。
不能用于别名、重算、spill或额外GET模型。

令原指令数为$N_i$、下界为$L_i$。逐图保留并同时报告$L_i/N_i$和$L_i/(N_i+L_i)$，二者均随COPY数单调增加。
原336图审计耗时5.460秒，其状态为269输出、47身份模型容量证书、12预算失败、8启发失败。
全部336图等权下界分别为25.960\%和6.690\%。仅269输出子集分别为31.128\%和7.278\%；子集不能替代全体评分。
初始使用输出分母时，全部图的2\%在该模型下不可能，而10\%尚未被该初始下界排除；后续更强必要界见下文。
47例模型容量证书还意味着在此身份模型下不能对这336图全部生成输出。
这些不是原题不可能性证明，因为官方未限定全部等价变换，也未提供实际评分数据分布。

\subsection{强制DAG区间及共同持续存活核心}
取严格祖先锚点$A\prec B$，选$k<W$个在$A$前出生、且有直接consumer严格在$B$后的原身份。
任何合法顺序均须使它们从$A$完成后持续存活至$B$之前。
其最新物理副本年龄互异，年龄和范围长度为$k(W-k)$。
任何不重新产生这些身份的原指令使年龄和增加$k$。
COPY可以引用较老副本；若引用年龄为$d_{\rm used}\le W$，最新副本年龄$d_{\rm latest}\le d_{\rm used}$，
刷新后的变化为$k-d_{\rm latest}\ge k-d_{\rm used}\ge-(W-k)$。
令$L$为$A$严格后代与$B$严格祖先交集的原节点数，则
\[
 C\ge\max\left(0,\left\lceil\frac{kL-k(W-k)}{W-k}\right\rceil\right).
\]
再选择一组共享GET身份，取所有严格在每个身份出生之后、又强制在每个身份未来直接消费之前的核心原节点。
这些节点可以互相不可比，但必在共同出生后到最早末消费前的区间内；以核心节点数代入同一年龄和不等式。
各证书区间可能重叠，只取最大值，绝不求和。
提议使用位集合；独立原JSON检查采用正、反DFS，不读取调度器闭包。
8192节点、1024子集和250毫秒是离线提议限制，超限保留下界0或已检查证书，不能解释为无解。

\subsection{多行共享归约的累加器角色下界}
严格完整结构有$R$行、$m$个共享权重GET、每行$m$个私有GET及MULW，原ADD链逐项归约到独立SET。
要求$R<W$、$m>W$。当前累加器身份随该行原ADD更新，但所有原节点及有序参数保持，证明没有授权结合重写。
最后一行首乘积出生必在第一SET之前：若某行尚未启动，每个已出生权重列必须保留该权重或该行未能归约的乘积，至多$W$列。
在两端点之间，$R$个累加器角色均存活；原GET/MULW增长年龄和$R$，原ADD或COPY最多减少$W-R$。
端点前后的GET/MULW分别至多$(2R+2)W$和$2RW$，因此
\[
 G=\max(0,(2R+1)m-(4R+2)W),\qquad
 C\ge\max\left(0,\left\lceil\frac{RG}{W-R}\right\rceil-R(m-1)-R\right).
\]
$R=4,W=8$时$C\ge\max(0,5m-144)$；4095原指令、315列的实例至少1431COPY，输出分母至少25.896\%。
3138不是因此最优，1431也不保证可达；此前必要界0不能说明能做到10\%。
独立原JSON证书逐行核对权重、私有GET、MULW、原ADD链、SET及全部节点覆盖，不调用调度器结构识别。

\subsection{新必要界审计}
\begin{center}
\begin{tabular}{lrrr}
\toprule
同一336输入的证据 & 全336 $L/(N+L)$ & 既有272合法 & 耗时\\
\midrule
R93强制锚点 & 11.2321\% & 9.5181\% & 7.988秒\\
R98共同核心 & 11.6440\% & 9.9195\% & 18.549秒\\
R99加共享归约 & 11.7792\% & 10.0864\% & 13.451秒\\
\bottomrule
\end{tabular}
\end{center}
完整336始终保留47个身份容量证书、9预算失败、8启发失败。R99原指令分母必要均值43.638\%，既有272子集35.267\%。
每阶段对272已完成独立身份检查的输出逐项核验下界不超过实际COPY，零矛盾；这只是排错，数学论证提供必要性。
独立raw年龄状态Dijkstra枚举全部就绪原节点和活身份COPY，没有使用ParsedBlock、PreparedGraph或求解器search。
R94新种子1200图：1057完整最优、3完整身份无解、140预算未知；25份正锚点证书中23与完整最优对照，零反例，耗时25.421秒。
R97专门正图192例、覆盖四窗口：168完整最优、24预算未知，零反例，耗时6.666秒。
新核心证书在R94出现96份（93完整最优、3完整无解），在R97出现192份（168完整最优、24未知），均无矛盾。
已加入COPY合法引用较老副本且达到窗口边界的独立原JSON测试，避免证明误假设输出总引用最新副本。

\begin{center}
\begin{tabular}{lrr}
\toprule
合成家族（原336图中各16例） & 平均$L/(N+L)$ & 局限 \\
\midrule
近窗口长脊柱 & 67.619\% & 大量接力具有结构必要性 \\
宽扇出 & 21.753\% & 包含8例身份容量不可行 \\
多输出不规则共享 & 10.110\% & 仅身份与COPY模型 \\
跨组共享 & 10.134\% & 包含证书及预算失败 \\
多输出固定共享 & 9.139\% & 下界不保证可达 \\
梯形菱形 & 7.568\% & 下界不保证可达 \\
\bottomrule
\end{tabular}
\end{center}

\section{新算法及复杂度}
\subsection{逆向压力候选}
从没有未调度后继的原节点开始逆向选择。
选择一个节点时，从逆向活跃集合删除其输出，加入有序输入所对应的不同身份。
优先不增加活跃集合的选择，再结合Sethi--Ullman需求、局部解锁量和活跃出生次序。
实现使用增量堆及缓存解锁计数，避免每次反复扫描高入度SET的所有先驱。
对每条数据和必要内存边更新有限次，排序与堆部分为$O((N+E)\log N)$，内存$O(N+E)$。
这不是COPY目标的完整算法；发射成本另外是$O(N+E+C)$，$C$为实际插入COPY数。
三个实验策略是压力、局部cluster近似及逆向跨度。
未声明任何一种在共享DAG上最优，固定候选失败也不是整体无解。

\subsection{严格共享根链}
独立结构筛选只接受：纯GET/SET/ADD/MULW，GET名字唯一、SET名字唯一且读写集合分离；
所有算术输出有一个不同consumer，沿SET回溯得到互不交叉的链；每条链的每个算术都使用同一个GET，
首次至多额外使用一个仅该处消费的私有GET，后续另一个输入来自前一链节点。
不同共享根组件之间没有数据关联；无用GET仍保留。
把共享根出生推迟到组件开始，再完成其分支，最后处理下一组件，可降低不必要长期活跃身份。
候选保持每个算术原有操作数顺序，GET不额外读取，SET不重复。
分支比较长先、短先、原顺序；只有COPY严格下降才物化新输出。
识别改用不同consumer计数，不构建无关通用优先级或需求图，复杂度$O(N+E)$。
模拟与物化均有时间上限，已经生成的合法候选在后续预算结束时保留。
新显式组合保留原有候选作回退；没有仅凭形状提前宣称COPY最优。

\subsection{梳状森林的早期零COPY构造}
完整纯森林中，每个非SET节点有一个不同consumer；若每个算术至多有一个算术前驱，其余输入是GET，
则从每个SET反向沿算术链回溯，正向执行时将私有GET紧贴消费点发射。
首次计算使用至多两个GET，后续计算使用前一步结果和至多一个新GET，最后SET紧贴结果。
重复操作数只读取同一原身份。全部物理引用距离至多2，因此四个官方窗口均合法，且COPY为0，达到非负下界。
筛选拒绝共享身份、无用GET、双算术前驱及寄存器hazard；不改变任何原算术、有序输入或架构访问。
识别、遍历和物化均为$O(N+E)$，避免通用PreparedGraph和需求排序。
最初放在容量预检之后只能救回3例；最终将这一构造性合法见证提前，避免在已有零COPY见证的结构上扫描无关的不可行证书。
候选仍经独立语义验证；外部原JSON身份重放证明它没有依赖等值别名。

\subsection{共享输入出生阶段和完整归约列发射}
\texttt{root\_phases}严格接受独立算术链共享GET的完整图，借用结构筛选但保持独立输出验证。
先推进只需要已经出生共享根的链段，再开启至多$W-S-1$条链，其中$S$为全图跨链共享GET身份数，留出物理接力余量。
可延迟新的共享GET出生；原GET各执行一次，原算术、有序参数和每个SET全部保留。
比较单链、双链、三链及较少后续使用的出生选择，仅在COPY下降时保留。
已知根集合单调增加，缓存各链可连续推进的末端，避免每步扫描完整剩余链。
每候选排序约$O((S+1)N+E)$，物化$O(N+E+C)$，最多四候选且有时间截止。
此限定邻域不证明全局最优，不能把不匹配形态或余量不足当成无解。
\texttt{phase\_portfolio}先保留旧候选再尝试这一邻域，默认\texttt{auto}仍保持旧策略。

\texttt{reduction\_columns}严格识别全覆盖共享权重乘积及原ADD链，逐列出生权重，再发射每行私有GET、MULW及对应原ADD。
还比较两个乘积一批、整列乘积一批及行次序变化，保持原二叉树，拒绝额外或遗漏节点。
$R=4$时$W\ge16$可零COPY，$W=8$的现有pressure已经产生相同列布局，不能称新质量收益。
两列26原节点$W=8$的固定序精确搜索完整22664状态，证明贪心8COPY在该顺序下最优，耗时202毫秒。
这只证明该顺序，不能排除不同列间重排有收益；因此没有以重复扩大同一固定序搜索当作新的进步。

\section{批次结果及负面证据}
所有批次保留全部输入、固定种子、原JSON哈希、策略成本、失败类型及成功输出；不删除困难图。
每个成功输出执行独立符号检查、原节点整数身份见证和三种数值差分。
以下速度是单进程逐输入配对，其他项目仍可能占CPU，有限预算可行数不具跨批绝对可比性。

\begin{longtable}{p{25mm}p{41mm}p{51mm}p{34mm}}
\toprule
批次 & 全部输入/成功 & COPY与失败 & 耗时 \\
\midrule
R82原336图 & v14 271；逆向压力248 & 共同248 COPY增加18531；17改善、65退化；丢失23输出 & 总106.661秒；不默认 \\
R83等长共享链40图 & v14 30；严格候选32 & 共同30 COPY9064相同；救回2；两者16384规模仍预算失败 & 总16.485秒；中位加速4.277倍 \\
R84新不等长共享根60图 & v14 52；严格候选60 & 共同52 COPY33927降到18650；4改善48相同；救回8，无退化 & 总28.306秒；中位加速3.149倍 \\
R85新长填充树32图 & v14及两候选各26 & 6例模型容量边界；原有work策略COPY9603相同，逆向压力7874 & 总8.130秒；详细失败另列 \\
R86共享根组合336图 & 旧/新各272 & 同47容量证书、9预算、8启发失败；COPY197152到195525，1改善271同 & 总68.931秒；中位0.990倍 \\
R87新组合336图 & 旧/新各272 & 12改善259同1退化；填充树4822到3902，非匹配图有预算波动 & 总67.862秒；中位0.993倍 \\
R88四种森林128图 & 旧97、新96 & 共同94：14改善80同；2救回3丢失，均在大梳状树压力发射阶段 & 总37.290秒；中位0.958倍 \\
R89梳状树48图初版 & 旧32、新35 & 零COPY但新13预算失败；容量预检仍消耗预算 & 改进仅显式，后续改早期见证 \\
R90同48图早期见证 & 旧29、新48 & 全新输出零COPY；19救回、共同29无质量变化 & 中位2.026倍；全成本另列 \\
R91完整336配对 & 旧/新各272 & 47身份证书、9预算、8启发失败；8改善264同0退化 & 总69.483秒；中位0.976倍 \\
R95有限消费者336图 & 冻结R91/新各272 & 3改善268同1退化；主要变化是旧候选预算，聚类本身无可观收益 & 中位0.981倍；仅显式 \\
R96共享新持出128图 & 旧108、新106 & 共同100一改善99同，6救回8丢失；唯一大幅COPY变化来自已有pressure完成 & 总62.732秒；不推荐 \\
R100共享根阶段336图 & 冻结R91/新各272 & 8改善264同；COPY194251到194178；全部73减少来自跨组结构候选 & 总58.188秒；中位0.987倍 \\
R101共享新持出64图 & 旧/新各60 & 同3证书1预算；21改善39同；COPY75048到74780，无输出丢失 & 总24.186秒；中位0.959倍 \\
\bottomrule
\end{longtable}

R82在全DAG上明显退化，说明寄存器压力目标与距离COPY目标不能等同。
其中一例COPY减少1933来自旧版本只能及时返回stable候选，不能称为超过旧affinity充分预算下的质量上限。
R83种子主要改变运算符，结构变化有限，所以新增R84的多根、偏斜长度、不同有序参数、首次重复参数与无用GET持出。
R84严格输出另从磁盘读回：60份符号及不同初值整数重放全部通过（5.959秒），原JSON身份重放60份全部通过（3.482秒）。
其全部60图等权实际COPY/输出为11.052\%、COPY/原指令为13.518\%；
独立下界分别为10.772\%和13.109\%（审计4.016秒）。
因此COPY/输出剩余空间至多约0.280个百分点，即使全部消除也达不到该批10\%。
旧成功子集52的均值与新60的均值不能当作相同样本直接比较。
R84的p95/p99仍为365/392毫秒；候选发射预算不包含无法中断的完整校验尾部，不能将250毫秒称作整个接口的严格墙时上限。
R85表明已有work-first森林策略并未取得新增收益；逆向单策略有16改善、6退化、4相同，所以组合先留旧森林候选，再尝试逆向。
R87非匹配近窗口图增加5COPY，且不规则共享图大幅收益主要来自旧有限预算没完成优质候选；不能把全部33066条减少都归因于算法质量。
R88的5个可行性差异均发生于原pressure候选执行阶段、逆向候选尚未运行，但失败仍保留，不以CPU竞争解释后删去。
R90采用R89相同输入，包含128、512、2048、8192、16384、32768节点，每种四窗口、两种新种子。
48份新输出全部经过四检查，旧29份可行，19份预算失败。全接口p95新604毫秒、旧461毫秒：新版本实际生成并完整验证了更大的输出，不能用所有输入的尾延迟直接声称加速。
共同29成功的中位加速2.026倍，零COPY提供局部Level比例的构造性证据，但不改变完整数据的Level不可达边界。

R91已完整通过141项回归（52.266秒）；10000固定种子小图两版本均9998合法、2身份模型证书，COPY均为562。
全部新输出通过符号、原JSON身份、三整数检查，失败保留依赖闭合子图范围的最小诊断反例；该子图可能删除SET，不能宣称其符合官方完整输入形态。
小图调度主循环约9.7秒，旧研究harness完整报告77.689秒，额外成本主要来自逐图扫描全部记录的二次配对，后续改为哈希索引并保留实际执行旧harness。
R91全336新旧同272输出，COPY195525到194736；固定272交集8改善264相同，没有质量或可行性退化。
新成功子集等权COPY/输出13.918\%、COPY/原42.007\%，分别仅是成功子集，不是全336均值。
以固定336分母计，已生成输出的非负贡献分别为11.267\%、34.006\%；其余64图无输出，因此完整成绩未定义。
不能把失败计为0伪称全体均值；若保持这272输出不变，任意非负补全也不能达到输出分母10\%。
272严格输出重新读回全部通过：符号与不同初值整数7.648秒，原JSON身份6.186秒。

R91基线实际来自冻结\texttt{experiments.baseline\_v14.isa\_scheduler.optimizer}，非当前核心的同名策略。
其规范化34文件聚合SHA为\url{e27b01946a217034074345db2546867ce38ebbeb41a78258dbd888b14dc41d1f}，与stage1冻结一致。
新R91核心SHA为\url{2543781e268a3f46fa2c75ff5689ad88d1668f4a4f819c3bb4d91e627aa6a7b2}。
R86、R87、R91相同源码和输入而COPY不同，是有限墙时下候选完成不同；只采用各轮同条件配对，不将跨轮的COPY差异笼统归因于历史算法改善。

R95/R96消费者聚类没有形成可归因的大幅收益，且新持出丢失8个旧成功输出，只保留显式实验和全部失败。
R100/R101的基线是实际\texttt{experiments/baseline\_level\_r91}冻结研究组合源码，而非当前核心同名策略。
R100仅跨组14个共同输出COPY9841到9768，目标1024节点$W=8$为1599到1573。
全272成功输出等权COPY/输出13.9100\%到13.8639\%，原指令分母41.9635\%到41.8751\%。
8个改善全部选择新根阶段策略，避免把其他策略因预算变化造成的质量差异误算作该候选收益。
R101新种子、128/512/2048/8192、四窗口、两家族，60共同成功等权COPY/输出19.9262\%到19.4634\%，原分母35.9302\%到34.9583\%。
质量收益实际可重复，但跨组中位速度分别0.853和0.822，时间成本增加；默认不替换旧算法。
R100严格输出272份重新读回原身份检查全部通过（3.922秒），符号及不同初值整数全部通过（4.918秒）。
R101严格输出60份也重新读回，原身份全部通过（2.002秒），符号及不同初值整数全部通过（2.618秒）。
R102阶段149项完整回归全部通过（49.798秒）；10000固定种子小图对实际冻结R91，两侧均9998合法、2身份模型证书，COPY均562。
新候选每份成功输出均经过原JSON身份、符号及三个整数配置检查，完整差分harness为7.657秒。
这批小图未获得COPY收益，正收益只来自已明确匹配的共享根链结构，不泛化为全部DAG改进。
下界证明的强化与上述真实73/268COPY减少是不同成果，必要界不计入优化收益。

\section{硬件、测试成本与开发余量}
当前使用Windows、CPython 3.12.14标准库，Ryzen 9 8945HX（16核32线程），物理内存约31.3GiB。
没有GPU、FPGA、STRAIGHT硬件或LLVM运行依赖。
主求解单进程，输出及证明保留会消耗$O(N+C)$空间；此前16384节点近窗口样例普通工作集约510MiB，显式下界停止约245MiB。
这是已测特例，不能作为所有搜索的峰值保证。
建议最低4核CPU、8GiB内存、Python 3.11以上、2GiB证据磁盘；建议8核以上、16--32GiB、10GiB磁盘用于多版本及重复实验。
这些是规划建议，未在最低硬件执行本报告批次。单图精确oracle只用于已声明规模的小图，状态上界随窗口和图规模指数增加，不能推荐对全大图无限搜索。

下一步应先用结构限定候选保持旧基线质量与所有失败分母，再扩大真实不同拓扑和规模。
一名开发者估计：完整规范适配及异常语义2--4工作日；限定结构候选及独立证明3--5日；
官方或真实应用负载采集、分布统计3--5日；C++热点原型与内存计量5--8日；
小图CP-SAT独立模型和交叉证明4--7日；完整性能与提交协议验收2--3日。
合计约19--32工作日，属于估算而非已投入工时。
规范适配、负载采集及C++原型可由不同人员并行；证明边界和基线数据冻结需要统一负责人。
若官方授权重算、GET重读或等值替换，应先建立新合法性及资源成本模型，再重新推导下界，估计另需3--6日。
没有必要通过GPU或付费算力扩大有限状态搜索；最明显的限制是合法性范围、负载分布和接近下界的结构，而非硬件不足。

\section{复现与证据保留}
旧\path{project18-research.tex/PDF}及\path{milestone-v14}原件不覆盖。
新候选代码包括\path{reverse_pressure.py}、\path{root_spines.py}、\path{root_phases.py}与\path{reduction_columns.py}。
冻结旧核心在\path{experiments/baseline_v14}。
实验目录包含\texttt{report.json}、\texttt{records.csv}、\texttt{inputs.jsonl}、严格输出及原节点排列。
R84输入SHA-256为
\url{b4ac53ff4309e6d0ac0d01f9db9b214255ef98e3b6cec6cc8f1ace8c7fdbbcc1}。
\path{research/level-stage/source-review.md}保存正文复核和一手链接。

R91完整回归、10000小图及原336验收已经保留；后续R93--R101均为新目录，不复写历史。
复现新阶段测试使用\texttt{scripts/level-stage.ps1}，可选\texttt{-Strategy phase\_portfolio -Baseline level\_r91}；
默认命令仍为冻结v14和研究组合对照。每次指定全新Round，以避免覆盖输入、输出或失败证据。
三份完整必要界推导分别位于\path{research/level-stage/forced-cut-proof.md}、\path{global-core-proof.md}及\path{shared-reduction-proof.md}。

另有根代理完全不导入本项目模块的原JSON审计，独立检查138份新证书（其中2份共享归约），
用独立距离解析、DFS、全图角色覆盖和大整数公式重算同272成功子集必要均值10.086440523\%。
数学推导审阅与反例搜寻仍是不同证据；有限输入审计不构成对任意输入的自动形式化证明。

R102代码冻结于2026-10-11 06:28:24 UTC，独立快照\texttt{experiments/baseline\_phase\_r102}。
规范化核心聚合SHA-256为
\url{9881efb6006e21fb81ad3ae6fd4cb904687090a87cd78c4108f9f0ed456f0b5d}。
\texttt{r102-phase-validation/source-provenance.json}保留逐文件及报告哈希。
R101和R102执行源与该快照完全一致；R100只有\texttt{\_\_main\_\_.py}命令行策略枚举在执行期间更新，算法文件未变，原开始时manifest原样保留。

\begin{thebibliography}{9}
\bibitem{sugita} Shu Sugita et al. A Sound and Complete Algorithm for Code Generation in Distance-Based ISA. CC 2023.
\url{https://doi.org/10.1145/3578360.3580263}.
作者全文：\url{https://www.rsg.ci.i.u-tokyo.ac.jp/lab/en/assets/files/cc2023-sugita.pdf}.
\bibitem{koizumi} Toru Koizumi et al. Compiling and Optimizing Real-world Programs for STRAIGHT ISA. ICCD 2021.
\url{https://www.rsg.ci.i.u-tokyo.ac.jp/lab/assets/files/iccd2021-koizumi.pdf}.
\bibitem{sethi} Ravi Sethi and J. D. Ullman. The Generation of Optimal Code for Arithmetic Expressions. JACM, 1970.
\url{https://doi.org/10.1145/321607.321620}.
\bibitem{chen} Gang Chen. Optimal and Heuristic Min-Reg Scheduling Algorithms for GPU Programs. 作者预印本v4，2023-11-21.
\url{https://arxiv.org/abs/2303.06855v4}.
\bibitem{bridge} Fan Yang et al. Register Bridging: A Lightweight Microarchitectural Approach for Suppressing Relay Move Instructions. ICCD 2025.
\url{https://doi.org/10.1109/ICCD65941.2025.00103}.
\url{https://pure.nitech.ac.jp/en/publications/register-bridging-a-lightweight-microarchitectural-approach-for-s/}.
\end{thebibliography}
\end{document}
